\paragraph{Practical implications.}
Communicative function predicts pruning safety more reliably than message content alone. Models with terse acknowledgments (Qwen, median 4 tokens) tolerate aggressive pruning; models with verbose, error-carrying acknowledgments (Llama, 59.8 tokens) benefit from pruning because removal eliminates an error-injection pathway. The mechanism divergence traces to alignment training: Qwen's tuning yields formulaic confirmations occupying zero bits of the task-relevant channel~\citep{clark_contributing_1989}, while Llama's RLHF produces verbose responses that consume channel capacity with noise~\citep{shannon_1948}. Latency gains are most pronounced in API-mediated deployments where fixed per-call costs dominate; post-hoc pruning is more practical than instructing agents to suppress emergent acknowledgments (Appendix~\ref{app:boundary_conditions}).

\paragraph{Information-theoretic interpretation.}
For Qwen, acknowledgment content and task outcome share near-zero mutual information ($H(Y \mid m_{\mathrm{ack}}) \approx H(Y)$; 5.7\% bigram novelty), placing acknowledgments outside the task-relevant channel. For Llama, acknowledgments carry positive mutual information with error introduction. This connects to \citet{rizvi_martel_communication_2026}: their bounds characterize when communication helps; our protocol identifies which messages consume capacity without contributing. The perturbation-pruning dissociation echoes classical results where local effects fail to predict global outcomes~\citep{lecun_obd_1990, braess_1968}.

\paragraph{Grounding or artifact?}
Unlike human dialogue where acknowledgments serve as grounding signals~\citep{clark_contributing_1989, clark_brennan_1991}, LLM-generated acknowledgments appear to be behavioral artifacts of instruction-tuning: removal does not degrade task accuracy, and their surface form correlates with RLHF training intensity rather than communicative necessity. Not all human communication primitives transfer functionally to LLM agents.

\paragraph{Limitations.}
Key limitations include single-seed causal evidence, greedy-only decoding, and domain scope restricted to closed-form QA tasks (GSM8K, StrategyQA, ARC, MATH); a HumanEval pilot ($N{=}50$) confirms code generation as a boundary (Appendix~\ref{app:humaneval-pilot}). The boundary conditions identified validate the framework's diagnostic ability to flag unsafe configurations. Three-agent pilots provide preliminary scalability evidence (Appendix~\ref{app:multiagent}); extended discussion in Appendix~\ref{app:extended_limitations}.
